\subsection{Rechenregeln der Laplace Transformation}

\textbf{Linearität}
Die Laplace-Transformation ist ein \textbf{linearer Operator}. Das bedeutet, dass die Transformation einer Linearkombination zweier Funktionen der Linearkombination der jeweiligen Einzeltransformationen entspricht. Sind $a, b \in \mathbb{R}$ Konstanten und $f(t)$ sowie $g(t)$ Laplace-transformierbare Funktionen, so gilt:

\begin{equation}
[\mathcal{L}(a f(t) + b g(t))](s) = a [\mathcal{L}f](s) + b [\mathcal{L}g](s)
\label{eq:laplace_linearitaet}
\end{equation}


Diese Eigenschaft vereinfacht die Berechnung von Laplace-Transformierten zusammengesetzter Funktionen, da man die einzelnen Bestandteile getrennt transformieren kann.

Gegeben sei folgendes Beispiel:
\begin{equation}
    [\mathcal{L}(3\exp(4t) + 7\exp(-2t))](s)
    = 3[\mathcal{L}(\exp(4t))](s) + 7[\mathcal{L}(\exp(-2t))](s)
    = \frac{3}{s - 4} + \frac{7}{s + 2}
    \label{eq:laplace_bsp_linear}
    \end{equation}
    \cite[S.~883]{buch1}
    
\textbf{Streckung}
\begin{equation}
    [\mathcal{L}(f(ct))](s) = \frac{1}{c}[\mathcal{L}f]\left(\frac{s}{c}\right), \quad c > 0
    \label{eq:laplace_streckung}
    \end{equation}

Diese Gleichung beschreibt die sogenannte \textbf{Streckungsregel} der Laplace-Transformation. Sie gibt an, wie sich die Transformation verändert, wenn die Zeitachse durch einen positiven konstanten Faktor $c$ gestaucht oder gestreckt wird. Wird in der Zeitfunktion $t$ durch $ct$ ersetzt, so entspricht das im Bildbereich einer Division des Laplace-Arguments $s$ durch $c$ sowie einer zusätzlichen Multiplikation mit $\frac{1}{c}$.

Beispiel:
\begin{equation}
    [\mathcal{L}(\cos(3t))](s) = \frac{1}{3}[\mathcal{L}(\cos t)]\left(\frac{s}{3}\right)
    = \frac{1}{3} \cdot \frac{\frac{s}{3}}{\left(\frac{s}{3}\right)^2 + 1}
    = \frac{s}{s^2 + 9}
    \label{eq:laplace_bsp_streckung}
    \end{equation}
    \cite[S.~884]{buch1}

Diese Gleichung zeigt ein konkretes Beispiel der Streckungsregel, angewendet auf $\cos(3t)$, unter Verwendung der bekannten Laplace-Transformierten von $\cos(t)$.

\textbf{Dämpfung}
\begin{equation}
    [\mathcal{L}(\exp(-at)f(t))](s) = [\mathcal{L}f](s + a), \quad a > 0
    \label{eq:laplace_exponentielle_modulation}
    \end{equation}

Diese Gleichung stellt die \textbf{Dämpfungsregel} der Laplace-Transformation dar. Sie besagt, dass die Multiplikation einer Zeitfunktion $f(t)$ mit einem exponentiellen Ausdruck $\exp(-at)$ im Zeitbereich dazu führt, dass die Laplace-Transformierte von $f(t)$ im Bildbereich um $a$ nach rechts verschoben wird, also $\mathcal{L}(f(t)) \to \mathcal{L}(f)(s + a)$. Diese Regel erlaubt es, die Transformierte zusammengesetzter Funktionen effizient zu bestimmen.


Beispiel:
\begin{equation}
    [\mathcal{L}(\exp(-t)\exp(t))](s) = [\mathcal{L}(\exp(t))](s + 1)
    = \frac{1}{(s + 1) - 1} = \frac{1}{s} = [\mathcal{L}(1)](s)
    \label{eq:laplace_exponential_beispiel}
    \end{equation}
    \cite[S.~884]{buch1}

\textbf{Ableitung}
Sei $f$ auf $[0, \infty[$ mindestens differenzierbar, wobei die Ableitung $f'$ auf jedem endlichen Intervall nur endlich viele Sprungstellen besitzt (also stückweise stetig ist). Alternativ kann $f$ auch stetig differenzierbar auf $[0, \infty[$ sein, dann kann folgende Rechenregel angewendet werden \cite[S.~23]{buch3}:

\begin{equation}
    [\mathcal{L}(f')](s) = s[\mathcal{L}f](s) - f(0)
    \label{eq:laplace_ableitung}
    \end{equation}


Beispiel:
\begin{equation}
    [\mathcal{L}(\sin' t)](s) = s[\mathcal{L}(\sin t)](s) - \sin(0)
    = s \cdot \frac{1}{s^2 + 1} = [\mathcal{L}(\cos t)](s)
    \label{eq:laplace_ableitung_beispiel}
    \end{equation}
    \cite[S.~884]{buch1}
    

Bei höheren Ableitungen muss iterativ vorgegangen werden, wobei die Anfangsbedingungen einbezogen werden
und mit steigender Potenz von $s$ der Grad der Ableitung fällt:

\begin{equation}
[\mathcal{L}(f^{(n)})](s)
= s^n [\mathcal{L}f](s) - f^{(n-1)}(0) - s f^{(n-2)}(0) - \dots - s^{n-1} f(0)
= s^n [\mathcal{L}f](s) - \sum_{k=0}^{n-1} s^k f^{(n-1-k)}(0)
\label{eq:laplace_hoehere_ableitung}
\end{equation}
\cite[S.~884]{buch1}

\textbf{Stammfunktion}

\begin{equation}
\left[ \mathcal{L} \left( \int_0^t f(u) \, du \right) \right](s)
= \frac{1}{s} [\mathcal{L}f](s)
\label{eq:laplace_integralregel}
\end{equation}

Diese Rechenregel beschreibt die Laplace-Transformation eines Integrals. Wenn man eine Funktion $f(t)$ zuvor über das Intervall $[0, t]$ integriert, so entspricht dies im Laplace-Bild einer Division der ursprünglichen Transformierten durch $s$.


Beispiel:

\begin{equation}
\left[ \mathcal{L} \left( \int_0^t \cos(u) \, du \right) \right](s)
= \frac{1}{s} [\mathcal{L}(\cos t)](s)
= \frac{1}{s} \cdot \frac{s}{s^2 + 1}
= [\mathcal{L}(\sin t)](s)
\label{eq:laplace_integral_beispiel}
\end{equation}
\cite[S.~885]{buch1}

\textbf{Faltung}

\begin{equation}
    [\mathcal{L}(f * g)](s) = [\mathcal{L}f](s) \cdot [\mathcal{L}g](s)
    \label{eq:laplace_faltung}
    \end{equation}
    
Die Faltungseigenschaft der Laplace-Transformation erlaubt es, das Faltungsintegral zweier Funktionen im Zeitbereich in eine einfache Multiplikation ihrer jeweiligen Laplace-Transformierten im Bildbereich umzuwandeln. Im gezeigten Beispiel wird die Faltung zweier identischer Exponentialfunktionen $e^t$ berechnet. Dabei ergibt sich, dass die Laplace-Transformierte des Faltungsprodukts $(e^t * e^t)(t)$ gleich der Laplace-Transformierten von $t e^t$ ist, was durch beide Rechenwege bestätigt wird.

Beispiel:
    
    \begin{align}
    [\mathcal{L}(e^t * e^t)](s)
    &= \left( [\mathcal{L}(e^t)](s) \right)^2
    = \left( \frac{1}{s - 1} \right)^2
    = \frac{1}{(s - 1)^2} \nonumber \\
    \text{Andererseits ist:} \quad
    e^t * e^t
    &= \int_0^t e^{t - u} e^u \, du
    = t e^t,
    \quad [\mathcal{L}(t e^t)](s)
    = \frac{1}{(s - 1)^2}
    \label{eq:laplace_faltung_beispiel}
    \end{align}
    \cite[S.~885]{buch1}
    

